\documentclass{book}
\usepackage[utf8]{inputenc}
\usepackage[T1]{fontenc}
\usepackage{geometry}
\usepackage{graphicx}
\usepackage{float}
\usepackage{hyperref}
\usepackage[english,indonesian]{babel}

\title{Optimasi Pemrosesan Template LaTeX Menggunakan Pendekatan Berbasis AI}
\author{Academic Support Agent}
\date{\today}

\begin{document}
\maketitle

\begin{abstract}

\end{abstract}

\chapter{PENDAHULUAN}


\chapter{TINJAUAN PUSTAKA}


% Original content (commented out):
% % --------------------------
% CREDITS:
% This template is based and basically an adaptation of UPM_thesis_template_en template:
% https://www.overleaf.com/latex/templates/upm-master-template-espanol/znqtbbyzjrnp
% All credits for the template to the author: MAría Blanco (UPM)
% --------------------------

% --------------------------
% Document class
% --------------------------
\documentclass[
letter,                % Tamaño del papel Carta
twoside,                % A una columna o a dos
openright,              % Los capítulos y secciones empiezan siempre a la derecha
chapterprefix=true,     % Añade un prefijo a los títulos
12pt,                   % Tamaño de fuente
headings=normal,        % Tamaño de las cabeceras
titlepage=on            % Separa las páginas con título
]{book}

% **************************************************
% THESIS details
% **************************************************

\newcommand{\University}{\protect{Universidad CENFOTEC}}
\newcommand{\UNIVERSITY}{\protect{UNIVERSIDAD CENFOTEC}}

\input{thesisDetails}

% **************************************************
% Settings
% **************************************************

\input{settings}

% **************************************************
% Begin document
% **************************************************
\begin{document}

% --------------------------
% Cover and front matter
% --------------------------
\frontmatter
\pagestyle{empty}				    % no header nor footer
\input{preliminares/1_cubierta}
\cleardoublepage

\pagenumbering{roman}			    % roman page numbering 
\pagestyle{plain}				    % empty header, just page number in footer
\input{preliminares/3_dedicatoria}
\cleardoublepage

\input{preliminares/4_agradecimientos_y_financiacion}   % dedication & aknowledgement (optional)
\clearpage


\input{preliminares/5_resumen}  % abstract (in English and Spanish)





\selectlanguage{spanish}

\setcounter{tocdepth}{3}		% define depth of toc
\tableofcontents				% display table of contents


\addcontentsline{toc}{section}{Índice de figuras}
\listoffigures


\addcontentsline{toc}{section}{Índice de tablas}
\listoftables
\clearpage

\input{preliminares/6_acronimos}
\clearpage


% --------------------------
% Main matter
% --------------------------
\mainmatter
\pagenumbering{arabic}			% arabic page numbering
\setcounter{page}{1}			% set page counter
\pagestyle{fancy} 	            % fancy header and footer
\parskip=7pt                    % space after each paragraph
\parindent=0pt                  % suppress indentation of paragraphs


\include{cuerpo/chapter1}
\include{cuerpo/chapter2}
\include{cuerpo/chapter3}
\include{cuerpo/chapter4}
\include{cuerpo/chapter5}
\include{cuerpo/chapter6}

\clearpage


% --------------------------
% Back matter
% --------------------------
\backmatter

\bibliographystyle{cas-model2-names}
\renewcommand\bibname{Referencias Bibliográficas}
\bibliography{references}


\appendix
\include{cuerpo/x-Anexos}


% **************************************************
% End of document
% **************************************************
\end{document}


\end{document}